%%%% main.tex -- HCAI 2026/27 project report (IJCAI--ECAI 26 format)
%%%% Draft. Every \todo{} marks content that depends on experiments not yet run
%%%% or on decisions still open; nothing inside a \todo{} is a real result.

\documentclass{article}
\pdfpagewidth=8.5in
\pdfpageheight=11in

\usepackage{ijcai26}

\usepackage{times}
\usepackage{soul}
\usepackage{url}
\usepackage[hidelinks]{hyperref}
\usepackage[utf8]{inputenc}
\usepackage[T1]{fontenc}
\usepackage[small]{caption}
\usepackage{graphicx}
\usepackage{amsmath}
\usepackage{amsthm}
\usepackage{booktabs}
\usepackage{multirow}
\usepackage{algorithm}
\usepackage{algorithmic}
\usepackage{xcolor}
\usepackage[switch]{lineno}

% Line numbers help reviewers; comment out for the final version.
\linenumbers

\urlstyle{same}

% Draft helpers: placeholders are printed in red so none is mistaken for a result.
\newcommand{\todo}[1]{\textcolor{red}{[TODO: #1]}}
\newcommand{\tbd}{\textcolor{red}{--}}   % empty table cell, value pending

\pdfinfo{
/TemplateVersion (IJCAI.2026.0)
}

\title{Human-in-the-Loop Active Learning for Breast Cancer Diagnosis\\from Ultrasound Images}

\author{
    Cristiana Mortágua \and Íris Sousa
    \affiliations
    Faculty of Sciences and Technology, University of Coimbra\\
    MSc in Artificial Intelligence and Data Engineering --- Human-Centered Artificial Intelligence
    \emails
    uc2023212004@student.uc.pt, uc2021221971@student.uc.pt
}

\begin{document}

\maketitle

% =====================================================================================
\begin{abstract}
Labelling medical images requires scarce expert time, and in breast cancer screening the
cost of an error is asymmetric: a missed malignancy (false negative) can delay treatment,
whereas a false alarm leads to further exams. We study whether a convolutional neural
network (CNN) trained with pool-based Active Learning (AL), in which a human annotator is
asked to label only the images the model considers most informative, can match a model
trained on the fully labelled set while keeping false negatives low. Using the public
Breast Ultrasound Images dataset (780 images; normal, benign, malignant), we first correct
two methodological problems of a previous implementation of this pipeline --- data leakage
caused by augmenting before splitting, and a Human-in-the-Loop step whose labels were
ignored by the learner --- and then compare three uncertainty-based query strategies
(least confidence, margin and entropy) under four training batch sizes against a passive
learning baseline. Evaluation focuses on the recall of the malignant class and on the
number of false negatives. \todo{one or two sentences with the main quantitative findings
once the experiments are run}. The code is available at
\url{https://github.com/httpirsh/miacd-iach-breast-cancer-diagnosis}.
\end{abstract}

% =====================================================================================
\section{Introduction}

Breast cancer is the most frequently diagnosed cancer in women worldwide and the leading
cause of cancer death among them~\cite{bray2024globocan}. In Portugal, about 9{,}000 new
cases are detected and more than 2{,}000 women die from the disease every year; around 1\%
of the cases occur in men~\cite{lpcc}. Survival depends strongly on how early the disease
is found, and ultrasound is a widely available, radiation-free imaging modality used both
to characterise lesions found in mammography and as a first-line exam in dense breasts and
younger women. Reading ultrasound images, however, is operator-dependent and demands
experience, which motivates computer-aided diagnosis tools.

Deep learning, and CNNs in particular, have reported high accuracies on breast ultrasound
classification~\cite{heliyon2023fusion,kaggle_anaya}. Two obstacles separate these
results from clinical use. First, CNNs need many \emph{labelled} images, and in medicine
each label costs the time of a specialist. Second, a model that classifies an image
without involving or informing the clinician is hard to trust, audit or correct: in a
high-stakes domain, the human must keep control of the decision~\cite{shneiderman2020hcai,holzinger2016interactive}.

Active Learning addresses the first obstacle: instead of labelling every image, the model
iteratively selects the unlabelled images whose labels it expects to be most informative
and asks an oracle --- here, a human expert --- to label them~\cite{settles2009survey}.
Because the expert is consulted inside the training loop, AL is also a natural instance of
Human-in-the-Loop (HitL) machine learning~\cite{mosqueira2023hitl,budd2021survey},
addressing part of the second obstacle. A previous student project on the same dataset
combined a CNN with AL through the modAL library and reported that AL reached 68\% test
accuracy against 45.92\% for passive learning~\cite{cardoso2024}. On inspection, that
pipeline had two problems that make its numbers unreliable: images were augmented
\emph{before} the train/test split, so rotated copies of test images were seen during
training (data leakage~\cite{kaufman2012leakage}); and in two of the three query
strategies the label typed by the human was not the one used to update the model, so the
``human'' in the loop had no effect. Moreover, accuracy alone was reported, although the
clinically relevant error is the false negative.

In this paper we revisit that study with a corrected and documented pipeline. We
(i)~split the data into stratified training, validation and test sets \emph{before} any
augmentation and remove a duplicated image with conflicting labels; (ii)~implement a
single AL loop in which the label given by the human --- real or simulated --- is the one
the model learns from, while the model's own suggestion is displayed to the annotator;
(iii)~compare least-confidence, margin and entropy sampling with four training batch
sizes against a passive learning baseline that uses the same CNN; and (iv)~evaluate
models by malignant-class recall and number of false negatives, alongside accuracy and
macro-averaged recall and F1. \todo{add (v) Grad-CAM explanations and (vi) the patient-facing
agent prototype once implemented, or move them to Future Work}.

Our approach connects directly to the topics of the Human-Centered AI course: it treats
the clinician as an active participant whose knowledge steers training (HitL rather than
fully automated learning~\cite{hcai_slides_hitl}), it chooses the evaluation criterion
according to human values (the asymmetric cost of missing a cancer), and it plans
explanations so that the human can judge \emph{why} the model asks for a label.

The remainder of this paper is structured as follows. Section~\ref{sec:related} reviews
related work. Section~\ref{sec:materials} describes the dataset and tools, and
Section~\ref{sec:methods} the pre-processing, model, AL loop and evaluation.
Section~\ref{sec:experiments} details the experimental setup, Section~\ref{sec:results}
presents the results and Section~\ref{sec:discussion} discusses them.
Section~\ref{sec:conclusions} concludes and outlines future work.

% =====================================================================================
\section{Related Work}
\label{sec:related}

\paragraph{CNN-based breast ultrasound classification.}
The Breast Ultrasound Images (BUSI) dataset~\cite{aldhabyani2020busi} is the most used
public benchmark for this task. Reported results vary widely: a CNN feature extractor
combined with an SVM reached 95.6\% accuracy on benign versus malignant
classification~\cite{heliyon2023fusion}, and public notebooks report accuracies above
99\%~\cite{kaggle_cbis}. Such figures should be read with care. BUSI is small, contains
duplicated images and images with visible annotations from the ultrasound
equipment~\cite{pawlowska2024curated}, and many pipelines augment before splitting. Under
these conditions a CNN can exploit shortcuts unrelated to the lesion~\cite{geirhos2020shortcut,zech2018confounding}, and leaked augmented copies inflate test scores. We
therefore do not aim to beat these numbers, but to measure AL against a passive baseline
under a protocol free of leakage.

\paragraph{Active Learning.}
In pool-based AL~\cite{settles2009survey}, a model trained on a small labelled set scores
a pool of unlabelled examples with a \emph{query strategy} and requests labels for the
best-scored ones. Uncertainty sampling~\cite{lewis1994sequential} queries the example on
which the model is least confident; margin sampling~\cite{scheffer2001margin} the one with
the smallest gap between its two most likely classes; and entropy
sampling~\cite{shannon1948} the one with the most spread-out predictive distribution.
For deep models, Bayesian approximations~\cite{gal2017deep} and diversity-based
strategies such as core-sets~\cite{sener2018coreset} have been proposed to deal with
over-confident softmax outputs and with batches of redundant queries.
Budd et al.~\cite{budd2021survey} survey AL and HitL for medical image analysis and
highlight that most studies simulate the oracle with ground-truth labels, as we also do
for the quantitative experiments. modAL~\cite{danka2018modal} provides these strategies in
Python, wraps any scikit-learn-compatible estimator and includes an example of
interactive labelling in Jupyter~\cite{modal_interactive}, on which our HitL interface is
based.

\paragraph{AL in medical diagnosis at the University of Coimbra.}
Pinto et al.~\cite{pinto2022active} combined modAL with a neural network for symptom-based
diagnosis of 41 diseases and obtained better performance than models in the literature
with less data. Faria et al.~\cite{faria2022joint} studied the joint role of batch size
and query strategy for heart attack prediction and found that small batches worked best
and that least-confidence, margin, entropy and random sampling all reduced the amount of
data needed by 57--60\%. Note that in~\cite{faria2022joint} ``batch size'' is the number
of examples queried per AL iteration, whereas in this work and in~\cite{cardoso2024} it
is the mini-batch size used by the optimiser when the CNN is (re)trained; one image is
queried per iteration. Vieira et al.~\cite{vieira2021openworld} applied AL with a
stacking ensemble to open-world object recognition for self-driving cars. Closest to our
work, Cardoso et al.~\cite{cardoso2024} applied the same CNN and the same three query
strategies to BUSI; we reuse their architecture to keep the comparison fair, and correct
the issues described in the introduction.

\paragraph{Human-centred and explainable AI.}
Shneiderman~\cite{shneiderman2020hcai} argues that high automation and high human control
are compatible and that reliable, safe and trustworthy systems keep humans in charge of
consequential decisions. Holzinger~\cite{holzinger2016interactive} discusses when a human
in the loop helps machine learning in health informatics, and Mosqueira-Rey et
al.~\cite{mosqueira2023hitl} review HitL paradigms, among which AL is the one in which the
machine drives the interaction. Showing the model's suggestion to the human can also
cause automation bias~\cite{goddard2012automation}, a risk we discuss in
Section~\ref{sec:discussion}. Grad-CAM~\cite{selvaraju2017gradcam} produces class-specific
heat maps from the gradients of the last convolutional layer and is commonly used to
check whether a CNN looks at the lesion or at artefacts.

% =====================================================================================
\section{Materials}
\label{sec:materials}

\subsection{Dataset}
We use the Breast Ultrasound Images dataset~\cite{aldhabyani2020busi}, distributed by
Dataset Ninja under a CC0~1.0 licence~\cite{datasetninja}. It contains 780 PNG images
from 600 women aged 25 to 75, acquired in 2018, labelled by radiologists as
\emph{normal} (133), \emph{benign} (437) or \emph{malignant} (210). Each image also comes
with a pixel mask of the lesion (absent for normal images), which we do not use for
training but which can serve as ground truth to evaluate explanations
(Section~\ref{sec:conclusions}). Image sizes vary between 190 and 1048 pixels in width
(median 564) and 310 and 719 pixels in height (median 474). The class of each image is
given by its file name (e.g.\ \texttt{benign (12).png}).

The dataset is imbalanced: there are more than three benign images for every normal one.
Furthermore, comparing file contents with MD5 hashes revealed that \texttt{benign (433)}
and \texttt{malignant (145)} are the same image with different labels. Since we cannot
know which label is correct, both copies are excluded, leaving 778 images (133 normal,
436 benign, 209 malignant).

\subsection{Tools}
The project is implemented in Python with TensorFlow/Keras for the CNN, scikit-learn for
data splitting and metrics, modAL~\cite{danka2018modal} for the AL loop, and SciKeras to
expose the Keras model through the scikit-learn interface that modAL expects; NumPy,
Pillow, Matplotlib and seaborn are used for data handling and plots. Shared code lives in
a \texttt{src/} package (data loading, model, AL loop, evaluation) and the experiments are
run from three Jupyter notebooks (data exploration, passive learning, active learning)
that work both locally and on Google Colab, which provides the GPU used for training.
\todo{state the hardware (Colab GPU type) and library versions used for the final runs}.

% =====================================================================================
\section{Methods}
\label{sec:methods}

\subsection{Pre-processing and data splits}
\label{sec:preprocessing}
Images are converted to a single grey-level channel, since ultrasound is intrinsically
monochromatic and three identical channels would only triple memory. They are resized to
$224\times224$ pixels, the input size used in~\cite{cardoso2024} and by most ImageNet
models (which keeps the door open to transfer learning), and pixel values are scaled to
$[0,1]$ to help optimisation. At this size the 1{,}899 images held in memory after
balancing take about 381\,MB as \texttt{float32}, well within the RAM of a free Colab
session. Resizing to a square changes the aspect ratio of the (mostly landscape)
images; we accept this distortion for simplicity.

The 778 original images are then split, stratified by class, into training (65\%),
validation (15\%) and test (20\%) sets. The split is done \textbf{before} any
augmentation, so no validation or test image, nor a rotated or flipped copy of it, is ever
used for training. The validation set is used to monitor training (and, in the AL loop,
is passed to every refit); the test set is used only for the reported evaluation. Both
keep the original images and the real class distribution, so that the reported metrics
reflect the prevalence the model would face.

\begin{table}[t]
\centering
\small
\begin{tabular}{lrrrr}
\toprule
Split & Normal & Benign & Malignant & Total \\
\midrule
Training (original)   & 86  & 283 & 136 & 505 \\
Training (balanced)   & 516 & 566 & 544 & 1626 \\
Validation            & 20  & 66  & 31  & 117 \\
Test                  & 27  & 87  & 42  & 156 \\
\bottomrule
\end{tabular}
\caption{Number of images per class in each split (seed 1). Only the training set is
balanced; the augmentation factors are $\times6$ (normal), $\times2$ (benign) and
$\times4$ (malignant).}
\label{tab:splits}
\end{table}

\paragraph{Balancing.}
Only the training set is balanced, by adding rotated and mirrored copies of its own
images: five extra copies of each normal image (rotations of 90$^\circ$ and 180$^\circ$,
a vertical flip, and both rotations combined with the flip), three of each malignant
image (180$^\circ$ rotation, flip, and both) and one of each benign image (flip). The
factors were chosen to bring the three classes to similar sizes (Table~\ref{tab:splits})
so that the model is not pushed to predict the majority class, which would hurt exactly
the minority classes we care about. The training set is then shuffled, because the
augmentation function appends copies class by class. We discuss the limitations of this
choice in Section~\ref{sec:discussion}.

\subsection{CNN model}
We keep the CNN of~\cite{cardoso2024}, which allows a direct comparison with that work and
isolates the effect of our methodological corrections. It has four convolutional blocks
with 32, 64, 128 and 32 filters of size $4\times4$, ReLU activation and no padding, each
followed by batch normalisation~\cite{ioffe2015batchnorm} and $2\times2$ max pooling;
dropout~\cite{srivastava2014dropout} with rate 0.3 is applied after the second and third
blocks. The resulting $11\times11\times32$ feature map is flattened and fed to a dense
layer of 32 sigmoid units and a 3-unit softmax output (Table~\ref{tab:cnn}). The model
has 355{,}203 parameters and is trained from scratch with Adam~\cite{kingma2015adam}
(learning rate $10^{-5}$) on the categorical cross-entropy loss, with one-hot labels. The
low learning rate favours stable training of a network that is retrained many times on a
growing set in the AL loop.

\begin{table}[t]
\centering
\small
\begin{tabular}{llr}
\toprule
Layer & Output shape & Parameters \\
\midrule
Input                         & $224\times224\times1$ & 0 \\
Conv $4\times4$, 32 + BN + Pool  & $110\times110\times32$ & 672 \\
Conv $4\times4$, 64 + BN + Pool + Dropout & $53\times53\times64$ & 33{,}088 \\
Conv $4\times4$, 128 + BN + Pool + Dropout & $25\times25\times128$ & 131{,}712 \\
Conv $4\times4$, 32 + BN + Pool  & $11\times11\times32$ & 65{,}696 \\
Flatten + Dense 32 (sigmoid)  & 32 & 123{,}936 \\
Dense 3 (softmax)             & 3  & 99 \\
\midrule
Total                         &    & 355{,}203 \\
\bottomrule
\end{tabular}
\caption{CNN architecture (\texttt{src/model.py}).}
\label{tab:cnn}
\end{table}

\subsection{Passive learning baseline}
The baseline trains a fresh CNN on the whole balanced training set (1{,}626 labelled
images) for 15 epochs with the default batch size of 32, monitoring the validation set
after each epoch. It represents the situation in which every training image has been
labelled by an expert, and is the reference that AL aims to approach with fewer labels.

\subsection{Active Learning with a human in the loop}
\label{sec:al}

\paragraph{Loop.} Algorithm~\ref{alg:al} summarises the procedure, implemented in
\texttt{src/active\_learning.py}. A random subset of $n_0$ training images forms the
initial labelled set $\mathcal{L}$; the remaining images form the unlabelled pool
$\mathcal{U}$. The CNN, wrapped by SciKeras with \texttt{warm\_start} enabled so that each
refit continues from the current weights, is first trained on $\mathcal{L}$. At each of
the $Q$ iterations, the query strategy scores every image in $\mathcal{U}$ and selects the
best one, the human provides its label, the image is moved from $\mathcal{U}$ to
$\mathcal{L}$, and the CNN is refitted on the whole of $\mathcal{L}$ for $E$ epochs
(modAL's \texttt{teach}). After every iteration we record loss and accuracy on
$\mathcal{L}$ and on the test set, and at the end we compute the test metrics.

\begin{algorithm}[t]
\caption{Pool-based Active Learning with HitL}
\label{alg:al}
\begin{algorithmic}[1]
\REQUIRE training set $X$, strategy $\phi$, $n_0$, $Q$, epochs $E$
\STATE $\mathcal{L} \leftarrow$ $n_0$ random labelled images; $\mathcal{U} \leftarrow X \setminus \mathcal{L}$
\STATE train CNN $f$ on $\mathcal{L}$ for $E$ epochs
\FOR{$t = 1$ \TO $Q$}
  \STATE $x^* \leftarrow \arg\max_{x \in \mathcal{U}} \phi(f(x))$
  \STATE show $x^*$ and suggestion $\arg\max_c f_c(x^*)$ to the human
  \STATE $y^* \leftarrow$ label given by the human
  \STATE $\mathcal{L} \leftarrow \mathcal{L} \cup \{(x^*, y^*)\}$; $\mathcal{U} \leftarrow \mathcal{U} \setminus \{x^*\}$
  \STATE continue training $f$ on $\mathcal{L}$ for $E$ epochs
  \STATE record loss and accuracy on $\mathcal{L}$ and on the test set
\ENDFOR
\end{algorithmic}
\end{algorithm}

\paragraph{Query strategies.} Let $p(c\mid x)$ be the softmax output for class $c$ and
$c_1, c_2$ the two most probable classes. We compare the three uncertainty-based
strategies available in modAL, which select the image maximising:
\begin{align}
\phi_{\mathrm{LC}}(x)  &= 1 - p(c_1\mid x) && \text{(least confidence)}\\
\phi_{\mathrm{M}}(x)   &= -\big(p(c_1\mid x) - p(c_2\mid x)\big) && \text{(margin)}\\
\phi_{\mathrm{H}}(x)   &= -\textstyle\sum_{c} p(c\mid x)\log p(c\mid x) && \text{(entropy)}
\end{align}
For example, for $p = (0.40, 0.35, 0.25)$, $\phi_{\mathrm{LC}} = 0.60$ and the margin is
$0.05$. With three classes the strategies often agree, but they can differ: least
confidence only looks at the top class, margin focuses on the boundary between the two
most likely classes (e.g.\ benign versus malignant), and entropy rewards images on which
the model hesitates among all three. Studying which of these notions of uncertainty finds
more useful images is the main question of the AL experiments.

\paragraph{The human in the loop.} When the \texttt{HUMAN} option is enabled, the selected
image is displayed in the notebook together with the model's suggested class, and the
annotator types the correct class; the model is trained with the label the person typed,
even if it is wrong. This is the behaviour a real deployment would have, and it fixes the
implementation of~\cite{cardoso2024}, in which the typed label was ignored for least
confidence and margin sampling. Running dozens of experiments with a person answering
every query is not feasible, so for the quantitative comparison the oracle is simulated
with the ground-truth label of the image, as is common in AL
studies~\cite{budd2021survey}; this corresponds to an always-correct expert.
\todo{if time allows, describe a small session with a real annotator (number of images,
agreement with ground truth, time per label)}.

\subsection{Evaluation metrics}
We report, on the test set: accuracy; recall and F1 macro-averaged over the three
classes; the recall (sensitivity) of the malignant class; and the number of false
negatives, defined as malignant images predicted as benign or normal. Accuracy is kept for
comparison with previous work but can be misleading with imbalanced classes. Malignant
recall and false negatives are the primary criteria, because in diagnosis a false positive
leads to additional exams, whereas a false negative can delay the treatment of a cancer.
Confusion matrices show how errors are distributed among classes.
\todo{add precision--recall curves and AUC-PR for the malignant class~\cite{saito2015precision}
if implemented}.

% =====================================================================================
\section{Experiments}
\label{sec:experiments}

\subsection{Experimental setup}
All experiments use the splits of Table~\ref{tab:splits} (seed 1), $224\times224$
images and the CNN of Table~\ref{tab:cnn}. The AL grid crosses the three query strategies
with four training batch sizes, $\{16, 32, 64, 128\}$, giving 12 runs; each run starts
from a fresh CNN. Following~\cite{cardoso2024}, each run starts with $n_0 = 1000$
randomly labelled images, the CNN is trained for $E = 10$ epochs at the start and after
every new label, and $Q = 10$ images are queried, one per iteration. The batch size is
varied because it changes how the optimiser exploits a small labelled set: small batches
give more, noisier updates per epoch, large ones fewer and smoother updates, and the
previous study found large batches to perform clearly worse~\cite{cardoso2024}. The
passive baseline is trained once as described in Section~\ref{sec:methods}.
Table~\ref{tab:params} lists all parameters; they are set in the first cells of the
notebooks.

\begin{table}[t]
\centering
\small
\begin{tabular}{ll}
\toprule
Parameter & Value \\
\midrule
Image size & $224\times224$, 1 channel \\
Split (train/val/test) & 65/15/20\%, stratified, seed 1 \\
Optimiser & Adam, learning rate $10^{-5}$ \\
Passive: epochs / batch size & 15 / 32 \\
AL: strategies & least confidence, margin, entropy \\
AL: training batch sizes & 16, 32, 64, 128 \\
AL: initial labelled set $n_0$ & 1000 (of 1626) \\
AL: iterations $Q$ (images queried) & 10, one per iteration \\
AL: epochs per (re)fit $E$ & 10 \\
Oracle & simulated (ground truth) \\
\bottomrule
\end{tabular}
\caption{Experimental parameters.}
\label{tab:params}
\end{table}

\todo{decide whether to add the control experiments proposed in
Section~\ref{sec:threats} (random sampling, compute-matched passive baseline, several
seeds, smaller $n_0$); if added, describe them here}.

% =====================================================================================
\section{Results}
\label{sec:results}

\todo{this section is a skeleton: fill in after running notebooks 01 and 02 with the
final configuration. Report numbers only; interpretation goes in
Section~\ref{sec:discussion}.}

\subsection{Passive learning baseline}
Table~\ref{tab:results} (first row) shows the test metrics of the passive baseline and
Figure~\ref{fig:cm} its confusion matrix. \todo{one or two sentences pointing out
accuracy, malignant recall and the number of false negatives, and whether training and
validation curves indicate under- or over-fitting after 15 epochs}.

\subsection{Active Learning}
Figure~\ref{fig:curves} shows the evolution of loss and accuracy on the labelled set and
on the test set across the 10 AL iterations, for each strategy and batch size, with the
accuracy of the passive baseline as reference. Table~\ref{tab:results} reports the test
metrics at the end of the 10 iterations. \todo{summary statements: which strategy/batch
size gave the fewest false negatives and the highest malignant recall; how many
configurations reached the baseline accuracy and after how many iterations; whether the
large batch sizes lagged behind}.

\begin{table}[t]
\centering
\small
\setlength{\tabcolsep}{4pt}
\begin{tabular}{llrrrrr}
\toprule
Strategy & Batch & Acc. & Rec. & F1 & Rec.\,mal. & FN \\
\midrule
Passive & 32 & \tbd & \tbd & \tbd & \tbd & \tbd \\
\midrule
\multirow{4}{*}{Least conf.} & 16  & \tbd & \tbd & \tbd & \tbd & \tbd \\
                             & 32  & \tbd & \tbd & \tbd & \tbd & \tbd \\
                             & 64  & \tbd & \tbd & \tbd & \tbd & \tbd \\
                             & 128 & \tbd & \tbd & \tbd & \tbd & \tbd \\
\midrule
\multirow{4}{*}{Margin}      & 16  & \tbd & \tbd & \tbd & \tbd & \tbd \\
                             & 32  & \tbd & \tbd & \tbd & \tbd & \tbd \\
                             & 64  & \tbd & \tbd & \tbd & \tbd & \tbd \\
                             & 128 & \tbd & \tbd & \tbd & \tbd & \tbd \\
\midrule
\multirow{4}{*}{Entropy}     & 16  & \tbd & \tbd & \tbd & \tbd & \tbd \\
                             & 32  & \tbd & \tbd & \tbd & \tbd & \tbd \\
                             & 64  & \tbd & \tbd & \tbd & \tbd & \tbd \\
                             & 128 & \tbd & \tbd & \tbd & \tbd & \tbd \\
\bottomrule
\end{tabular}
\caption{Test-set results (156 images, 42 malignant). Rec.\ and F1 are macro-averaged;
Rec.\,mal.\ is the recall of the malignant class; FN is the number of malignant images
predicted as benign or normal (lower is better). AL rows are measured after 10 queries
($1010$ labelled images).}
\label{tab:results}
\end{table}

\begin{figure*}[t]
\centering
\fbox{\parbox[c][4cm][c]{0.95\textwidth}{\centering\todo{three panels (one per strategy)
exported from notebook 02, \S4.1: accuracy and loss per iteration on the labelled set and
on the test set, one colour per batch size, passive baseline as a dashed line}}}
% \includegraphics[width=\textwidth]{figures/al_curves.pdf}
\caption{Accuracy (solid) and loss (dashed) per Active Learning iteration on the labelled
set (left) and test set (right) for each training batch size. The black dashed line is
the test accuracy of passive learning.}
\label{fig:curves}
\end{figure*}

\begin{figure}[t]
\centering
\fbox{\parbox[c][3.5cm][c]{0.9\columnwidth}{\centering\todo{confusion matrices: passive
baseline and best AL configuration}}}
% \includegraphics[width=\columnwidth]{figures/confusion.pdf}
\caption{Confusion matrices on the test set. Rows are true classes and columns predicted
classes; false negatives are in the \emph{malignant} row, \emph{normal} and
\emph{benign} columns.}
\label{fig:cm}
\end{figure}

% =====================================================================================
\section{Discussion}
\label{sec:discussion}

\todo{rewrite with the real results. The paragraphs below state what we expect and how
each outcome should be interpreted, so that the discussion is planned before the numbers
are known.}

\paragraph{Expected comparison with the passive baseline.}
After 10 iterations an AL model has seen 1{,}010 labelled images, 62\% of the 1{,}626
available to the baseline. If its malignant recall and false negatives are comparable to
the baseline's, AL would have saved the labelling of roughly 600 images; if they are
better, the explanation is not necessarily the query strategy (see
Section~\ref{sec:threats}). Because only 10 images are queried, we expect the
differences between strategies to be small and possibly within run-to-run variance.

\paragraph{Expected effect of the strategies.}
Margin sampling concentrates on the boundary between the two most probable classes. In
this dataset the hardest boundary is clinically the most important one --- benign versus
malignant lesions, which differ mainly in shape and margins --- so we expect margin to
query more benign/malignant images and, possibly, to reduce false negatives the most.
Entropy may also select normal images on which the model is uncertain among all three
classes. \todo{confirm or refute with the class distribution of the queried images}.

\paragraph{Expected effect of the batch size.}
With a learning rate of $10^{-5}$ and a fixed number of epochs, a batch of 128 performs
only about 8 weight updates per epoch on 1{,}000 images, against 63 for a batch of 16. We
therefore expect large batches to train more slowly and stay further from the baseline
within the 10 iterations, as observed in~\cite{cardoso2024}; a better learning rate or
more epochs would be needed for a fair comparison of batch sizes.

\paragraph{The role of the human.}
In the quantitative experiments the oracle never errs. A real annotator can be tired,
inexperienced, or influenced by the model's suggestion, which we display to make labelling
faster: if the suggestion is wrong, the annotator may accept it (automation
bias~\cite{goddard2012automation}), and the model would then be reinforced in its own
mistake, since it learns from the human's answer. This trade-off between efficiency and
independence of the human judgement is central to HitL design; hiding the suggestion, or
showing it only after the human answers, are alternatives worth testing. Explanations
(Grad-CAM) would let the expert judge whether the model's uncertainty comes from the
lesion or from artefacts such as calipers and text burned into the image.

\subsection{Threats to validity}
\label{sec:threats}
We list the limitations of the current protocol, which also motivate the control
experiments we plan.

\emph{Unequal training budget.} Each AL refit trains for 10 epochs on the whole labelled
set, starting from the previous weights. After 10 queries, an AL model has processed about
$10\times1000 + 10\times\sum_{k=1}^{10}(1000+k) \approx 1.1\times10^5$ training examples,
against $15\times1626 \approx 2.4\times10^4$ for the passive baseline, i.e.\ about $4.5$
times more gradient steps. Any advantage of AL over the baseline may therefore come from
longer training rather than from the choice of images. A passive baseline trained for the
same number of steps, or with early stopping on the validation set, removes this
confound.

\emph{No random-sampling control.} Without a run in which the queried images are chosen at
random~\cite{faria2022joint}, we cannot attribute differences to the uncertainty
criteria themselves.

\emph{Large initial set and few queries.} Starting from 1{,}000 of 1{,}626 images (61.5\%)
and querying 10, the labelled set grows by only 1\%; AL is most useful when the initial
set is small and many queries are made. A smaller $n_0$ (e.g.\ 100) with more queries
would test the label-efficiency claim more directly.

\emph{Augmented copies in the pool.} The pool contains rotated and flipped copies of
training images, possibly of images already labelled. A real expert would label an image
once, and augmentation would be applied afterwards; querying a copy of a known image costs
a label but brings little information.

\emph{Augmentations that are not physically plausible.} Ultrasound images have a fixed
orientation (skin at the top, acoustic shadows below lesions); 90$^\circ$ and 180$^\circ$
rotations produce images a scanner never outputs. Horizontal flips, small rotations,
speckle-like noise and contrast changes, applied on-the-fly, would be more realistic.

\emph{Single run.} Results come from one data split and one initialisation; differences of
one or two false negatives in a 42-image malignant test set are within noise. Repeating
each configuration with several seeds and reporting means and standard deviations would
allow statistical comparison.

\emph{Model selection.} Choosing the ``best'' strategy by its test-set false negatives
uses the test set for selection; the validation set should be used instead, with the test
set reserved for the final report.

% =====================================================================================
\section{Conclusions and Future Work}
\label{sec:conclusions}

\todo{"looking back": summarise what was achieved and the main quantitative finding}.

We revisited a CNN with Active Learning and a human in the loop for breast ultrasound
classification. The main contributions so far are methodological: a pipeline that splits
the data before augmentation and removes a mislabelled duplicate, an AL loop in which the
human's answer actually drives learning, and an evaluation centred on false negatives.
The main difficulties were \todo{e.g.\ training time on Colab, memory, the previous code
base}.

\paragraph{Relation to Human-Centered AI.}
The project touches three course topics. \emph{Human-in-the-Loop}: the expert is consulted
only where the model is uncertain, keeping human effort low while leaving the final
authority over labels with the human. \emph{Values and evaluation}: optimising for false
negatives reflects the patient's interest rather than a generic accuracy. \emph{Trust and
explainability}: explanations and an honest description of limitations are prerequisites
for a clinician to rely on, or overrule, the system~\cite{shneiderman2020hcai}.

\paragraph{Future work.}
(i)~Grad-CAM explanations~\cite{selvaraju2017gradcam} shown to the annotator with each
query, evaluated against the lesion masks of the dataset with IoU or Dice;
(ii)~a patient-facing agent prototype, in a notebook, that presents the prediction and its
explanation and answers questions about the disease through retrieval-augmented
generation~\cite{lewis2020rag} over trusted sources (e.g.~\cite{lpcc}), with a clear
statement that it is not a medical diagnosis; (iii)~transfer learning from ResNet or
EfficientNet~\cite{he2016resnet,tan2019efficientnet}; (iv)~the controls of
Section~\ref{sec:threats}, plus diversity-based and committee-based query
strategies~\cite{sener2018coreset,settles2009survey}; and (v)~a user study with
clinicians to measure labelling time, agreement and the effect of showing the model's
suggestion. \todo{move items to the main text as they get implemented}.

% =====================================================================================
\section*{Reproducibility}
The repository (\url{https://github.com/httpirsh/miacd-iach-breast-cancer-diagnosis})
contains the code, the notebooks, a \texttt{requirements.txt} and a README with
instructions to obtain the dataset and run the notebooks in order, locally or on Google
Colab; all parameters are set in the first cells of each notebook and metrics are saved
to \texttt{results/}. The \LaTeX{} sources of this paper are in \texttt{report/}.

%% \section*{Contribution Statement}
%% \todo{who did what, if required}

\bibliographystyle{named}
\bibliography{references}

\end{document}
